\chapter*{Conclusion and perspectives}
\addcontentsline{toc}{chapter}{Conclusion and perspectives}
\markboth{Conclusion and perspectives}{}

\paragraph{}This final-year project focused on the design and implementation of a configurable test data
generator for Linedata Longview. The initial problem came from the difficulty of preparing coherent
test datasets manually for a financial platform whose data model contains many dependent entities:
accounts, account groups, hierarchies, models, securities, prices, positions, orders, and
executions. Manual SQL preparation was time-consuming, error-prone, difficult to reproduce, and
limited by confidentiality constraints that prevent the direct use of production data.

\paragraph{}The developed solution addresses this problem through a full-stack web application composed of a
React frontend and a NestJS backend. The system generates Longview-compatible pre-BCP records,
validates them, and assembles them into a SQL loading script that can be downloaded by QA engineers.
The standalone mode allows fully synthetic dataset
generation without a database connection, while the connected mode allows generated records to be
anchored to existing SQL Server accounts or securities selected by the user.

\paragraph{}The project was carried out using an Agile Scrum organization divided into five sprints. Sprint 1
mapped the Longview pre-BCP buffer model and established the population order. Sprint 2 delivered
the core generation engine and standalone workflow. Sprint 3 introduced SQL Server connectivity and
connected generation. Sprint 4 added optional LLM enrichment to improve the realism of descriptive
financial fields while preserving algorithmic fallback. Sprint 5 finalized the configuration
library, allowing users to save, load, rename, and delete reusable generation scenarios.

\paragraph{}Several engineering choices contributed to the reliability of the final tool. The generator was
implemented as a set of entity-specific services, making the business rules easier to isolate and
maintain. Validation is performed before file writing and SQL assembly, reducing the risk of
delivering invalid datasets. Connected mode uses read-only SQL queries and restricts raw table
fetching to approved Longview tables. LLM integration is optional and defensive: it enriches only
selected descriptive fields and never replaces deterministic generation as the source of structural
truth.

\paragraph{}The final application therefore meets the main objectives of the internship. It reduces manual
SQL preparation effort, improves reproducibility through saved configurations, supports both
standalone and connected test scenarios, and produces a coherent SQL script aligned with the
Longview loading process.

\paragraph{}Several perspectives remain possible for future work. The first improvement would be to extend
the generator to additional Longview entities and client-specific custom fields. A second direction
would be to add automated execution of the generated SQL script against a dedicated test database,
including a post-load verification report. A third perspective would be to introduce scenario
templates for common QA cases, such as full order lifecycle testing, fixed income schedule testing,
or derivative instrument testing. Finally, broader testing across multiple Longview environments
would help validate connected mode against different database configurations and client-specific
schema variations.
